\scanpage{591}
\mspage{599}
\uline{Remark}\quad The elements of the inverse \unsure{image|image}
$S_0\Gamma_{\rho,\sigma}^{\unsure{'!|!'}}$ of $Z_{\rho,\sigma}^{!'}$
in $S_0\Gamma_{\rho,\sigma}^{!}$ which belong to the section under consideration
are characterized (by virtue of the proposition on page 598) by the relation
\[
(61)\qquad u\equiv b\equiv a\quad\bmod\widehat\Pi_0
\]
\margin{(in addition to $\varepsilon_3(u)=1$ which characterizes the elements of
$\unsure{S_0\Gamma'_{\rho,\sigma}}$)}
which makes it evident that this is a subgroup.
I have not found a nice analogous expression for the membership of an
element of $S_0\Gamma_{\rho,\sigma}^{!''}$
(characterized by the relation $\varepsilon_2(u)\varepsilon_3(u)=1$)
in the corresponding section.

\uline{Proposition}\quad Let
$S_0^\Pi\Gamma_{\rho,\sigma}^{!}\subset S_0\Gamma_{\rho,\sigma}^{!}$
be the subset of $S_0\Gamma_{\rho,\sigma}^{!}$
\deleted{\unsure{of}}
\[
(62)\quad
\begin{gathered}
S_0^\Pi\Gamma_{\rho,\sigma}^{!}\subset S_0\Gamma_{\rho,\sigma}^{!}\\
\unsure{''}\\[-1ex]
\left\{(u,a,b)\in
\begin{array}{c}
\inserted{Z_{\rho,\sigma}\times N_\rho^*\times N_\sigma^*}\\[-.5ex]
\deleted{\unsure{S_0\Gamma_{\rho,\sigma}^{!}}}\\[-.5ex]
\inserted{M[0]}
\end{array}
\ \middle|\ 
\begin{array}{l}
u\equiv b\quad(\widehat\Pi_0)\\
u\equiv\deleted{\unsure{b\equiv}}\varepsilon_3(\sigma)a\quad
\unsure{\bmod}\ \widehat\Pi=\widehat\Pi_0\times\{\pm1\}
\end{array}
\right\}
\end{gathered}
\]
Then

$a)$ $S_0^\Pi\Gamma_{\rho,\sigma}^{!}$ is a subgroup of
$S_0\Gamma_{\rho,\sigma}^{!}$

$b)$ The \unsure{hom\textsuperscript{omorphism}} \deleted{\unsure{is}}
\diagram{\begin{tikzcd}[ampersand replacement=\&,row sep=small,column sep=large]
S_0^\Pi\Gamma_{\rho,\sigma}^{!}\arrow[r]
\& Z_{\rho,\sigma}^{!}=M^{!}[0],\\
(u,a,b)\arrow[r,mapsto]\&u
\end{tikzcd}}
is surjective, and its kernel is the subgroup, isomorphic to
$\mu=\{\pm1\}$, of elements $(1,\varepsilon,1)$ with
$\varepsilon\in\mu=\deleted{\unsure{S_0\Gamma_{\rho}}}
N'_\rho\cap\widehat\Pi=L_\rho\cap\widehat\Pi$
\mspage{590}

\scanpage{592}

\uline{Proof} The fact that we have a subgroup follows from the calculation already made on pages\ $597',598$, the condition that the product of two elements $(u',a',b')(u,a,b)$ of $S_0^{\Pi}\Gamma_{\rho,\sigma}$ lie in it still being expressed by the formula
\[
\tag{56}
[\dot\sigma,\dot u']=\dot\varepsilon'_3
\qquad\text{\uline{in}}\quad M/\widehat\Pi
\quad(\text{instead of }M/\widehat\Pi_0)
\]
\deleted{\unsure{We have}}\inserted{or again}
\[
\tag{56 bis}
\begin{gathered}
[\dot u',\dot\sigma]=\dot\varepsilon'_3
\qquad\text{i.e.}\quad \dot\varepsilon'_3=\dot\varepsilon'_2\\
\unsure{\text{''}}\\
\dot u'(\dot\sigma)\dot\sigma^{-1}=\dot\varepsilon'_2
\end{gathered}
\]
Now in $\unsure{\widetilde{\mathcal G}|\widetilde{\mathcal C}}/\Pi_0$, \unsure{the} images of $\varepsilon'_2$ and $\varepsilon'_3$ both belong to the central subgroup $\Pi/\Pi_0$, hence they are trivial in $\unsure{\widetilde{\mathcal G}|\widetilde{\mathcal C}}/\Pi$, hence relation (56 bis) is \unsure{trivially|trivially} satisfied.

The surjectivity of the map $S_0^\Pi\Gamma^!_{\rho,\sigma}\longrightarrow Z^!_{\rho,\sigma}$ comes from the fact that $S_0^\Pi\Gamma^!_{\rho,\sigma}$ contains the set-theoretic section $A(Z^!_{\rho,\sigma})$ of the proposition on page\ 597. The kernel consists of the \deleted{$\unsure{u}$} $(1,a,b)$ with $a\in N_\rho^*\cap\Pi$, $b\in N_\sigma^*\cap\Pi_0$.

Now $N_\rho^*\cap\unsure{\widetilde{\mathcal G}|\widetilde{\mathcal C}}=\deleted{\unsure{Z_\rho}}L_\rho$, $N_\sigma^*\cap\unsure{\widetilde{\mathcal G}|\widetilde{\mathcal C}}=L_\sigma$, and
\[
\deleted{L_\rho\cap\unsure{\Pi_0|\Pi}=\unsure{\mu}\,L_\sigma\cap\unsure{\Pi}}
\]
\[
\tag{57}
L_\rho\cap\widehat\Pi=L_\sigma\cap\widehat\Pi=\mu,
\qquad L_\rho\cap\widehat\Pi_0=L_\sigma\cap\widehat\Pi_0=\{1\}
\]
\mspage{591}

\scanpage{593}
\mspage{600}

whence the characterization of the kernel.

\uline{Corollary} Let us consider \deleted{$\unsure{Z}$} $S_0\Gamma^!_{\rho,\sigma}$ as an extension of $Z^!_{\rho,\sigma}=M^![0]$ by
\[
\deleted{\unsure{\varepsilon|Z}}\quad
L_\rho\times L_\sigma\simeq\mathbb Z/6\mathbb Z\times\mathbb Z/4\mathbb Z,
\qquad\text{hence}
\]
(writing $L_\rho\simeq\inserted{\mu\times}L_{-\rho}\deleted{\times\unsure{\mu}}\simeq\mathbb Z/3\mathbb Z\times\mathbb Z/2\mathbb Z$)
an extension of
\[
\underbrace{\mathbb Z/2\mathbb Z}_{\mu}\times
\underbrace{\mathbb Z/3\mathbb Z}_{L_{-\rho}}\times
\underbrace{\mathbb Z/4\mathbb Z}_{L_\sigma}.
\]

\deleted{The extension \unsure{co}}

This extension \unsure{corresponds} therefore \unsure{canonically|canonically} to three extensions, by $\mu,L_{-\rho},L_\sigma$ respectively. \unsure{The|The} \inserted{\unsure{extensions|extension}}\deleted{sections} by $L_{-\rho}\times L_\sigma$ \deleted{are} equipped \unsure{two|with a} canonical sections.

There remains an extension (this is precisely $S_0^\Pi\Gamma^!_{\rho,\sigma}$) of $Z^{\unsure{!}}_{\rho,\sigma}$ by $\mu$.

This extension corresponds \unsure{just} to the ambiguity of the sign in choosing an $\alpha=\widehat\Pi$ over $\widetilde\alpha\in\widehat\Pi/\pm1$ defined by $u\in Z^!_{\rho,\sigma}=M^![0]$, which is reflected in the ambiguity in choosing the sign for $a\in\unsure{N'_\rho|N_\rho^*}$, linked to $\alpha$ by (51)
\[
\tag{58}
a=\varepsilon_3(\tau')\alpha^{-1}u,
\qquad \alpha=ua^{-1}\varepsilon_3(\tau'),
\]
We were able to remove this ambiguity by the requirement
$\alpha\in\widehat\Pi_{\unsure{0\tau|0,\tau}}=\deleted{\unsure{\Pi_0}}\{1,\tau\}.\widehat\Pi_0$, and we

\mspage{592}

\scanpage{594}

have just observed that this choice, canonical though it is, is not multiplicative. At the same time, we find the complete explicitation of the $\alpha''$ (\unsure{normalized} \unsure{canonically|canonically} by the condition $\alpha''\in\widehat\Pi_{0\tau}$) associated with $u''=u'u$, in terms of $\alpha'$ (associated with $u'$) and $\alpha$ (associated with $u$) also normalized, \inserted{in terms of $a',a$} Indeed, the proposition on page\ 598 tells us in which cases \deleted{\unsure{exact|exactly}} we have
\[
\tag{59} a''=a'a
\]
namely iff we have
\[
\tag{59 bis}
\varepsilon_3(u)=1\quad\text{or}\quad
\varepsilon_{\deleted{\unsure{3|2}}\inserted{2}}(u')\varepsilon_3(u')=1.
\]
But we now know that in any case we must have
\[
\tag{60}
a''=\varepsilon a'a\qquad \varepsilon\in\{\pm1\}
\]
and we now find that $\varepsilon=+1$ iff we are in case (59 bis), $-1$ in the contrary case i.e. if $\varepsilon_3(u)=-1$ and $\varepsilon_2(u')\varepsilon_3(u')=+1$. In summary

\uline{Corollary} Let $u,u'\in M^![0]$, whence $u''=u'u$, let $\alpha,\alpha',\alpha''\in\widehat\Pi_{0\tau}$ be associated (defined by $u(\rho)=\mathrm{int}(\alpha)\rho$ etc) and put

\mspage{593}

\scanpage{595}
\mspage{601}

\[
\tag{61}
a=\varepsilon_3(\tau')\alpha^{-1}u,\qquad
a'=\varepsilon'_3(\tau')\alpha'^{-1}u',\qquad
a''=\varepsilon''_3(\tau')\alpha''^{-1}u''
\]
where
\[
\varepsilon_3(\tau')=
\begin{cases}
1&\text{if }\varepsilon_3=1\\
\tau'&\text{if }\varepsilon_3=-1
\end{cases}
\qquad
\text{and the }\overline{\mathrm m}\text{ for }\varepsilon'_3(\tau'),\ \varepsilon''_3(\tau'),
\]
and where $\varepsilon_3=\varepsilon_3(u)=\varepsilon_3(\unsure{\chi|\gamma}(u))$,
and likewise for $\varepsilon'_3,\varepsilon''_3$
--- hence $\varepsilon''_3=\varepsilon_3\varepsilon'_3$. We then have
\[
\tag{62}
a''=\varepsilon a'a,\quad\text{with}\quad\varepsilon\in\{\pm1\}
\]
and
\margin{\deleted{\unsure{and}}\ \unsure{therefore}}
\[
\tag{63}
\begin{cases}
\varepsilon=+1&\text{iff}\quad \varepsilon_3=\unsure{1}\ \text{\uline{or}}\quad
\varepsilon'_2\varepsilon'_3=1\qquad\text{hence}\\
\varepsilon=-1&\text{iff}\quad \varepsilon_3=\unsure{-1}\ \text{and}\quad
\varepsilon'_2\varepsilon'_3=-1.
\end{cases}
\]

We \unsure{will have|have} therefore
\deleted{(64)}
\[
\tag{\inserted{*}}
\left\{
\begin{aligned}
\alpha''&=u''a''^{-1}\varepsilon''_3(\tau')\\
&=\varepsilon u'u\,a^{-1}\underbrace{a'^{-1}\varepsilon_3^{\unsure{\deleted{''}|''}}(\tau')}_{\vphantom{A}}\varepsilon'_3(\tau')
\end{aligned}
\right.
\]
I \deleted{\unsure{say|would say} $=\varepsilon u'(\alpha)$}
\inserted{\unsure{presume}} that we will have
\margin{not true in general}
\[
\tag{64}
\alpha''=\varepsilon u'(\alpha)\alpha'
\qquad\text{where }\varepsilon\text{ given by (63).}
\]

Indeed
\[
\left\{
\begin{aligned}
\alpha&=u\,a^{-1}\varepsilon_3(\tau')\\
\alpha'&=u'\,a'^{-1}\varepsilon'_3(\tau')
\end{aligned}
\right.
\qquad\text{whence}
\]
\[
\begin{aligned}
u'(\alpha)&=u'u\,a^{-1}\varepsilon_3(\tau')u'^{-1}\\
u'(\alpha)\alpha'&=u'u\,a^{-1}
\underbrace{\varepsilon_3(\tau')\deleted{u'^{-1}}\deleted{u'}a'^{-1}}_{\vphantom{A}}
\varepsilon'_3(\tau')
\end{aligned}
\]
Formula (64), by virtue of (*), thus takes the form
\[
a'^{-1}\varepsilon_3(\tau')=\varepsilon_3(\tau')a'^{-1}
\]
i.e.
\[
[\varepsilon_3(\tau'),a'^{-1}]=1
\]
it is not always satisfied if $\varepsilon_3\deleted{(\unsure{u})}=-1$
i.e.\ $\varepsilon_3(\tau')=\tau'$, where it represents \inserted{again} a
draconian restriction on $a'$. The formula that is always correct is
\deleted{\unsure{(64)|(65)}}
\[
\alpha''=(\varepsilon u'(\alpha)\alpha')\mathbin{.}
\underbrace{\varepsilon'_3(\tau')\bigl([a',\varepsilon_3(\tau')]\bigr)}_{[\varepsilon'_3(\tau')(a'),\,\varepsilon_3(\tau')]}
\]
\mspage{594}